\section{Does a VLA Know Where It Is?}

A VLA can recognize the task object while still failing to locate the body that
will act on it. We diagnose this missing-self failure with a two-stage probe.
First, we expose a behavioral symptom: when the physical manipulation state is
unchanged but the wrist-camera view is shifted, the predicted action can shift
with the image. Second, we inspect the model's visual reasoning and find that
robot-centric queries do not reliably localize the robot body itself. Together,
these probes suggest that image-space shortcuts arise because current VLAs can
lack visual self-grounding.

\subsection{VLAs Can Rely on Image-Space Shortcuts}

We first isolate whether the predicted action is anchored to the physical
robot-scene state or to the image frame. The robot proprioception, object
configuration, language instruction, and ground-truth action are held fixed,
while the wrist-camera view is laterally shifted to the left or to the right.
This perturbation changes where the scene appears in the image, but it does
not change the correct robot-centric action.

A spatially grounded policy should remain stable under this view-only
perturbation. Instead, we observe co-directional action drift: left-shifted
views induce left-biased trajectories, and right-shifted views induce
right-biased trajectories relative to the ground truth. This drift
substantially reduces task success, indicating that the policy can complete
training-view tasks without fully anchoring action prediction to robot-scene
geometry.

This behavior reveals an image-space shortcut. The policy can partially tie
control to pixel coordinates, camera-specific correlations, or object-location
priors learned from the training view. Under the original camera, these cues
can be predictive enough to support imitation. Under a lateral camera shift,
however, the same cues move in image space while the correct action remains
fixed, causing the action prediction to move with the view. We quantify this
effect with Trajectory Shift Ratio (TSR), which measures the magnitude of
prediction drift, and Shift Direction Agreement (SDA), which measures whether
the drift follows the image shift.

\subsection{The Missing Self-Grounding Capability in VLAs}

We next ask why image-space shortcuts remain available by probing whether the
model visually grounds the robot itself. In both simulation and real-robot
scenes, we query the VLM or VLA backbone with robot-centric prompts such as
``Where is the robot itself?'' and visualize the resulting attention or token
relevance maps. The prompt is only a diagnostic tool: because it explicitly
names the robot, a visually self-grounded model should assign high relevance
to the visible arm, gripper, or end-effector.

The attention maps often fail to concentrate on the robot body. Instead, they
assign relevance to task objects, tabletop regions, or other visually salient
context, even when the robot is clearly visible. This indicates that task
grounding and self-grounding are not the same capability. A VLA may identify
the instructed object and still fail to represent where the acting body appears
in the image.

This model-level probe supports our attribution of the behavioral failure.
Attention maps do not fully explain the policy's internal computation, but
they provide consistent evidence that current VLAs can lack visual
self-grounding. When the visual pathway does not encode where the robot body
appears, the policy can satisfy action supervision through image-space
correlations under the training camera. Those correlations then become brittle
shortcuts under lateral viewpoint shifts. Robust VLA control therefore needs a
training signal that binds proprioceptive state to visual evidence of the
robot body, rather than assuming that self-grounding emerges from language and
action supervision alone.
